\begin{abstract}
We characterize the precision of population average treatment-effect estimation in balanced randomized trials with selectively observed binary outcomes, a binary surrogate, and a finite baseline alphabet. Outcomes are missing at random conditional on treatment, baseline category, and surrogate, with unknown cell-specific arrival probabilities bounded below by a supplied positive floor \(q\). This floor indexes and tunes the procedure. For \(n\), the sample size, and \(d\), the number of baseline categories, we give an explicit mathematical attainment construction: a deterministic mixed-count estimator with squared-error risk bounded by a universal constant times the comparison rate
\[
r_{n,d,q}
=
\min\left\{1,\frac{1}{nq}
+\left[\frac{d}{nq\log(e+nq)}\right]^2\right\}.
\]
This guarantee holds for every positive arrival floor. Under \(q\log^2(e+nq)\le1/64\), matching lower bounds establish minimax squared-error order \(r_{n,d,q}\) and optimal maximal expected length of order \(\sqrt{r_{n,d,q}}\) for uniformly honest connected intervals at each fixed coverage level strictly between zero and one. At each fixed \(0<q<1/2\), the point-risk comparison also matches, with constants depending on \(q\). Near complete ascertainment, an observed-outcome contrast attains risk at most \(4/n+(1-q)^2\). For \(d\ge1024n^2\), the matching minimax order is \(n^{-1}+(1-q)^2\), making \(1-q=O(n^{-1/2})\) necessary and sufficient along arrival-floor sequences for parametric risk uniformly over all alphabet sizes. A target-preserving synthetic assignment transfers the construction to discrete observational regression contrasts and closes the stated logarithmic minimax gap at each fixed strict-interior treatment-positivity margin.
\end{abstract}

\section{Introduction}

The precision of a randomized trial at a prespecified analysis horizon depends on both the primary outcomes that have arrived and the information recorded for the full sample. Baseline categories, treatment assignments, and short-term measurements can remain available while primary outcomes are being collected. When outcome arrival varies across these observed characteristics, estimating a population treatment effect requires combining population cell weights with outcome information within cells. This paper characterizes the resulting statistical precision in a balanced randomized experiment with binary primary outcomes, a binary surrogate, and a finite baseline alphabet.

The model imposes independent and identically distributed sampling, treatment assignment independent of baseline characteristics and potential outcomes and surrogates, assignment probability one half, and consistency. Outcome arrival is missing at random conditional on treatment, baseline category, and realized surrogate. Each occupied cell has an unknown arrival probability at least a supplied value \(q\), where \(0<q\le1\). This lower bound indexes the model and is used to tune the estimator and interval; individual cell probabilities remain unknown. Baseline probabilities may vary arbitrarily, including zero probabilities. The surrogate supplies observed conditioning information for ascertainment, and the target is the population average treatment effect on the primary outcome. Under these conditions, \cref{obj:lem:unrestricted-cell-identification} expresses the effect as a signed sum of cell membership probabilities multiplied by outcome means among arrivals. Membership information from every sampled individual enters this representation.

Our first result is a constructive squared-error guarantee across the entire positive-arrival domain. Write \(n\) for the sample size, \(d\) for the number of baseline categories, and \(N=nq\) for the effective sample-size scale. The mixed-count estimator in \cref{obj:def:mixed-count-estimator} satisfies a uniform risk bound of order
\[
r_{n,d,q}
=
\min\left\{1,\frac1N+
\left[\frac{d}{N\log(e+N)}\right]^2\right\},
\]
where \(r_{n,d,q}\) is the comparison rate, with a universal comparison constant, as established in \cref{obj:thm:unrestricted-mixed-count-upper}. The construction separates membership, pilot, and outcome estimation into three auxiliary streams. Membership counts estimate cell weights; the pilot selects between empirical arrived-outcome means and a Chebyshev polynomial statistic built from marked factorial counts. Exact conditional averaging over the auxiliary randomization produces a deterministic estimator.

Under the rare-arrival restriction
\[
q\log^2(e+nq)\le\frac1{64},
\]
\cref{obj:thm:matched-minimax-frontier} establishes matching minimax squared-error order \(r_{n,d,q}\), with universal constants. The decision class includes all bounded parameter-independent randomized estimators. Its minimax value agrees with that for deterministic estimators by \cref{obj:lem:randomized-kernel-barycenter}. The lower bound retains the full observed experiment, including memberships whose outcomes have yet to arrive. A one-category family supplies the reciprocal effective-sample-size contribution, while an activated moment-matching construction supplies the category contribution. Along sequences satisfying the rare-arrival restriction and \(N\to\infty\), \cref{obj:prop:phase-boundaries} consequently gives consistency exactly when \(d=o\{N\log(e+N)\}\), and squared-error order \(N^{-1}\) exactly when \(d=O\{\sqrt N\log(e+N)\}\).

The same rare-arrival experiment admits a matching uncertainty-quantification comparison. For each fixed miscoverage level \(\alpha\in(0,1)\), \cref{obj:thm:connected-interval-frontier} establishes optimal maximal expected length of order \(\sqrt{r_{n,d,q}}\) among connected intervals with uniform coverage at least \(1-\alpha\). The comparison constants may depend on \(\alpha\). Coverage and expected length integrate jointly over the sample and any endpoint randomization. A deterministic squared-error envelope supplies an attaining interval, and the converse covers the full randomized endpoint decision class.

Near complete ascertainment, the remaining deficit supplies a complementary precision scale. Throughout the arrival-floor model, \cref{obj:thm:unrestricted-near-complete-arrival-envelope} gives
\[
\frac{1}{256n}
\le \mathfrak R^{+}_{n,d,q}
\le
\min\left\{1,\overline C r_{n,d,q},
\frac4n+(1-q)^2\right\},
\]
where \(\mathfrak R^{+}_{n,d,q}\) is the minimax squared-error risk and \(\overline C\) is a universal constant. The observed-outcome arm contrast attains the last upper bound. For \(d\ge1024n^2\), \cref{obj:thm:unrestricted-large-alphabet-deficit-lower} establishes matching order \(n^{-1}+(1-q)^2\) uniformly in \(q\). Hence, along arrival-floor sequences, an eventual \(n^{-1}\) risk guarantee holding simultaneously over every alphabet size is equivalent to \(1-q_n=O(n^{-1/2})\), by \cref{obj:prop:unrestricted-uniform-near-complete-elbow}. Complete arrival gives \(n^{-1}\) order uniformly over arbitrary alphabet growth.

The construction also connects randomized outcome arrival to observational adjustment. An independent fair assignment retains an observational outcome when the synthetic assignment agrees with the observed treatment. Under a fixed strict-interior treatment-positivity margin, this transformation preserves the baseline-weighted regression contrast and supplies a randomized missing-at-random completion. Averaging the mixed-count estimator over synthetic assignments yields the deterministic observational procedure in \cref{obj:thm:zeng-fixed-positivity-polynomial-frontier}. At each fixed margin strictly between zero and one half, its minimax squared-error order is
\[
\min\left\{1,\frac1n+
\left[\frac{d}{n\log(e+n)}\right]^2\right\},
\]
with margin-dependent constants. Consistency and conditional exchangeability give the observational contrast its population causal interpretation. The same experiment comparison supplies the matched randomized point-risk order at every fixed \(0<q<1/2\) in \cref{obj:thm:unrestricted-fixed-q-minimax-frontier}, with constants depending on that floor.

The statistical mechanism is a weighted-cell-mean problem with many scarce arrived counts. Ordinary cell ratios accumulate a category-dependent bias; on the difficult large-alphabet branch, a Chebyshev correction reduces the aggregate bias to the scale \(d/\{N\log(e+N)\}\). The converse uses moment-matched reciprocal priors whose activated observation laws remain close even though all cell memberships are retained. Each augmented record carries an independent activation flag with probability \(qH\), where \(H\) is the upper endpoint of the latent prior support and has squared-logarithmic order. This explains the rare-arrival condition ensuring a valid same-experiment construction. \Cref{sec:constructive-estimator,obj:thm:full-experiment-lower} present these two sides.

These results connect three strands of existing work. Surrogate-assisted analysis explains how early measurements inform treatment-effect estimation and efficiency under appropriate observation and nuisance conditions \citep[Theorems 2.1--2.2 and 4.1--4.2]{KallusMao2025}. Decaying-labelling analysis establishes asymptotic inference under effective-sample-size and nuisance-rate conditions \citep[Theorems 2.2 and 4.1, equation (2.3)]{ZhangChakraborttyBradic2025}. Our comparisons characterize uniform squared error and honest expected interval length as category frequencies, alphabet size, and ascertainment vary over the specified finite models. For discrete observational adjustment, \citet[Assumption 2, Theorems 1--2, Section 3.2, and Appendix C.5]{ZengBalakrishnanHanKennedy2026} provide the closest rate benchmark. The synthetic-assignment construction supplies an attaining polynomial estimator, while their lower-bound conclusion supports the fixed-positivity comparisons. The polynomial and moment-matching methods build on discrete functional estimation \citep{JiaoVenkatHanWeissman2015,Wu2016,WuYang2019}; the expected-length criterion connects to honest confidence-interval theory \citep{Donoho1994,Cai2004,ArmstrongKolesar2020}.



The exposition proceeds through related work, the experiment and identification conditions, and the point-risk comparisons across ascertainment regimes. It then develops the constructive estimator, honest intervals under rare arrival, and the observational extension through synthetic assignment. Discussion and open questions follow. The appendices collect identification and decision-theory preliminaries, risk certificates, full-experiment lower bounds, interval and observational arguments, the verification note, and proofs of the main results.
